\documentclass[11pt, a4paper]{article}

% --- Dados do documento ---
\newcommand{\docTitle}{Prompts para o Dia a Dia}
\newcommand{\docSubtitle}{Guia Prático de Engenharia de Prompts}
\newcommand{\docAuthor}{Fernando Anselmo}
\newcommand{\docCategory}{Material de Estudo e Prática}
\newcommand{\docVersion}{Versão 1.0}
\newcommand{\docDate}{\today}

\usepackage{../sty/tutorial2}

% --- Início do Documento ---
\begin{document}
	
% =====================================================
% CABEÇALHO DE TÍTULO
% =====================================================
\begingroup
\setlength{\parskip}{0pt}

% Linha decorativa superior (sem espaço extra antes)
{\color{secondary}\rule{\textwidth}{2pt}}

\vspace{0.3cm}

% Título principal
{\Huge\sffamily\bfseries\color{primary} \docTitle\par}

\vspace{0.2cm}

% Subtítulo
{\Large\sffamily\color{secondary} \docSubtitle\par}

\vspace{0.3cm}

% Linha decorativa inferior
{\color{secondary}\rule{\textwidth}{1.5pt}}

\endgroup

\textbf{Autor:} \docAuthor

\textbf{Categoria:} \docCategory

\textbf{Versão:} \docVersion

% =====================================================
% SEÇÃO INTRODUTÓRIA: SOBRE ESTA APOSTILA
% =====================================================
\textbf{Resumo:} Reúne conceitos, coletânea de exemplos e boas práticas sobre a \textbf{engenharia de prompts}, com o objetivo de capacitar o leitor a se comunicar de forma clara e eficiente com ferramentas de inteligência artificial generativa. Ao longo do material, você encontrará desde os fundamentos básicos até técnicas avançadas de refinamento, sempre com exemplos aplicáveis ao dia a dia.

\textbf{A quem se destina:} profissionais, estudantes e curiosos que desejam extrair o máximo de ferramentas como ChatGPT, Gemini, Claude e similares.

\textbf{Como usar:} leia os capítulos na ordem apresentada para uma progressão natural, ou consulte diretamente os tópicos de seu interesse conforme a necessidade.

% =====================================================
% SEÇÃO INTRODUTÓRIA: O QUE É ENGENHARIA DE PROMPTS
% =====================================================
\section*{Introdução}
\addcontentsline{toc}{section}{Introdução}

Antes de começarmos deixe-me explicar que nas seções dedicadas aos \textit{prompts} não pretendo declamar autoria sobre nenhum deles, o que fiz foi recolher, as vezes melhorar, pedir auxílio para a IA se o prompt está bom e capturar e guardar em uma Probola.

A engenharia de prompts é a prática de criar, testar e aprimorar instruções fornecidas às ferramentas de inteligência artificial (IA) para obter respostas mais precisas, relevantes e alinhadas aos objetivos desejados. Embora o termo prompt também seja utilizado na programação para indicar que um sistema aguarda uma interação do usuário, no contexto da inteligência artificial ele representa as instruções que orientam o modelo na execução de uma tarefa.

Neste contexto, atuarei como seu especialista em engenharia de prompts. Você poderá apresentar uma ideia, um problema ou um objetivo que deseja alcançar utilizando inteligência artificial. A partir dessas informações, farei perguntas para compreender suas necessidades, esclarecer os requisitos e identificar oportunidades de aprimoramento. Com base nas suas respostas, construirei um prompt claro, estruturado e otimizado, aplicando boas práticas de engenharia de prompts para aumentar a qualidade, a consistência e a relevância dos resultados obtidos. 

\subsection*{Melhore seu prompt}
O padrão: Contexto + Papel + Tarefa + Formato é uma excelente estrutura para criar prompts mais eficazes e obter respostas alinhadas aos seus objetivos. Assim sendo, a partir de agora, utilizaremos esses quatro elementos como referência para desenvolver nossos prompts:
\begin{itemize}
	\item \textbf{Contexto}: compreender o cenário, o problema e as informações necessárias para realizar a tarefa.
	\item \textbf{Papel}: definir a função, a especialidade ou a perspectiva que a inteligência artificial deverá assumir.
	\item \textbf{Tarefa}: especificar claramente o que deverá ser feito, incluindo objetivos e requisitos.
	\item \textbf{Formato}: estabelecer como a resposta deverá ser apresentada, considerando estrutura, extensão, linguagem e nível de detalhamento.
\end{itemize}

\textbf{Exemplo de prompt fraco}: "Escreve um e-mail para o cliente sobre o atraso na entrega." 

\textbf{Exemplo de prompt forte}: "Aja como um gerente de atendimento de uma empresa de móveis planejados. Escreva um e-mail para um cliente avisando que a entrega vai atrasar 5 dias por problema no fornecedor. O tom deve ser transparente, sem parecer que estamos nos desculpando demais, e o e-mail deve terminar oferecendo um desconto de 5\% como compensação. Máximo de 120 palavras." 

Facilmente observamos que o segundo prompt gera um resultado praticamente pronto para ser enviado, enquanto o primeiro exige uma boa dose de ajustes. A diferença está na forma como as instruções são elaboradas.

Essa é a diferença entre:
\begin{itemize}
	\item Prompt ingênuo: pede o que você quer e recebe o que pediu, mesmo que seja impossível.
	\item Prompt robusto: pede o que você quer e uma avaliação honesta de se aquilo é alcançável.
\end{itemize}

Um prompt forte não é necessariamente robusto, vamos realizar uma análise e correção desse prompt (E-mail de atraso)

É um prompt curto e claro, sendo que a maioria dos problemas dele é de calibragem, não de estrutura. Sete pontos:
\begin{enumerate}
	\item "Transparente, sem parecer que estamos nos desculpando demais" é uma instrução de tom, não de conteúdo. O modelo entende "não se desculpar demais" de formas diferentes: pode cortar o pedido de desculpas por completo, pode substituir por linguagem corporativa fria, ou pode manter o pedido mas sem repetição. Sem critério, o resultado é imprevisível. Falta dizer o que deve estar presente (reconhecimento do impacto, responsabilidade, ação concreta) e o que deve estar ausente (repetição de "desculpe", autoflagelação, justificativa longa).
	\item "Problema no fornecedor" é vago e coloca a culpa fora. Isso pode soar como terceirização de responsabilidade ("a culpa não é nossa"). Para o cliente, o fornecedor é invisível — ele comprou da sua empresa. O prompt deveria orientar como mencionar a causa sem transferir culpa e sem expor detalhe interno desnecessário.
	\item "Desconto de 5\% como compensação" sem contexto. 5\% de quê? Do valor total? Do saldo? É aplicado como abatimento, cashback, crédito? O modelo vai inventar. E "compensação" sem como e quando aplicar gera ambiguidade que vira novo problema no atendimento.
	\item Falta o insumo essencial: quem é o cliente e qual é o pedido. O prompt não inclui: nome do cliente, número do pedido, valor, data original, nova data. O modelo vai usar placeholders genéricos ([Nome], [pedido]) — o que é aceitável, mas precisa ser explicitamente pedido, senão ele inventa dados fictícios que parecem reais.
	\item "Máximo de 120 palavras" é bom, mas não diz o que priorizar. Em 120 palavras, você não cabe: saudação calorosa + explicação detalhada + pedido de desculpas elaborado + desconto + encerramento. O modelo vai cortar algo — e provavelmente vai cortar o pedido de desculpas ou a explicação, não o que você queria cortar. Falta dizer o que é essencial e o que é descartável.
	\item Falta o canal e o contexto do relacionamento. Um e-mail para cliente novo é diferente de um para cliente recorrente; um para cliente que já reclamou é diferente de um para cliente que ainda não sabe do atraso. O tom "transparente, sem excesso de desculpas" muda conforme o histórico.
	\item Falta instrução sobre o que não dizer. Em comunicação de crise, o que se evita é tão importante quanto o que se inclui: não prometer o que não se controla, não dar detalhes internos, não usar jargão corporativo, não minimizar o impacto ("apenas 5 dias"), não prometer compensação além do autorizado.
\end{enumerate}

\textbf{Há ainda um ponto silencioso}: o e-mail tem um objetivo implícito que o prompt não nomeia — preservar o relacionamento e evitar que o cliente cancele ou reclame publicamente. Sem esse objetivo explícito, o modelo escreve um e-mail "correto" mas que pode não atingir o que você realmente quer.

Agora vamos criar uma versão corrigida e realmente robusta dele:

\begin{promptbox}[E-mail de atraso de entrega ao cliente]
Aja como gerente de atendimento de uma empresa de móveis planejados, escrevendo em nome da empresa. Vou te dar os dados e as regras; escreva o e-mail final.

Dados do caso (se eu não preencher algum, use um placeholder entre colchetes, não invente valor):
 Cliente: [nome]
 Número do pedido: [número]
 Data original de entrega: [data]
 Nova data de entrega: [data]
 Causa: problema no fornecedor de [componente/insumo]
 Compensação autorizada: desconto de 5\% sobre [valor total / saldo restante], aplicado como [abatimento na próxima fatura / crédito / estorno]
 Histórico do cliente: [primeira compra / recorrente / já reclamou do atraso / ainda não sabe do atraso]

O que o e-mail precisa conter (essencial, nesta ordem):
 Assunto claro, com número do pedido (sem alarmismo, sem "URGENTE")
 Reconhecimento direto do atraso e da nova data - na primeira ou segunda frase
 Uma frase de responsabilidade da empresa (sem culpar o fornecedor publicamente; se mencionar a causa, diga "problema na nossa cadeia de fornecimento", não "o fornecedor falhou")
 Ação concreta: o que a empresa está fazendo para garantir a nova data
 A compensação de 5%, com como e quando será aplicada
 Um canal direto de contato para o cliente (nome + telefone/e-mail do gerente)
 
O que o e-mail NÃO deve conter (regras de exclusão):
 Repetição de "desculpe" (uma vez basta; o reconhecimento vem dos fatos, não da repetição)
 Autoflagelação ("falhamos vergonhosamente", "não temos desculpa")
 Justificativa longa sobre o fornecedor (o cliente não precisa do detalhe interno)
 Minimização do impacto ("apenas 5 dias", "pequeno atraso")
 Promessa que a empresa não controla ("garantimos que não vai atrasar de novo")
 Jargão corporativo ("sinergia", "alinhamento", "tratativa")
 Compensação além dos 5% autorizados
 Tom frio ou formal demais (é atendimento, não comunicado jurídico)

Tom:
 Transparente e direto: o cliente deve entender o que aconteceu e o que vai acontecer em uma leitura.
 Responsável sem subserviência: a empresa assume o problema e resolve.
 Humano, não corporativo: frases curtas, voz ativa, sem "informamos que".
 
Formato e limite:
 Máximo 120 palavras no corpo (assunto não conta).
 Se precisar cortar, corte explicação da causa primeiro, depois agradecimento, nunca o reconhecimento, a nova data ou a compensação.
 
Depois do e-mail, liste em uma linha: 
 O que ficou implícito (o que o cliente vai inferir) e o que ficou fora de propósito (o que você decidiu não incluir e por quê).

Regras gerais:
 Se algum dado essencial estiver em branco, use placeholder - não invente nome, valor, data nem número de pedido.
 Se você não tiver certeza sobre convenção de atendimento no setor de móveis planejados, sinalize.
 Sem introdução antes do e-mail; comece direto pelo assunto.
\end{promptbox}

As ferramentas de inteligência artificial generativa são capazes de realizar diversas tarefas, mas precisam de comandos claros e bem estruturados para oferecer resultados mais precisos e alinhados às expectativas do usuário. Quando um prompt é muito vago, as respostas tendem a ser genéricas ou diferentes do que se esperava. Por isso, quanto mais detalhes relevantes forem fornecidos, maiores serão as chances de obter um resultado satisfatório.

A engenharia de prompts é importante para aprimorar a interação entre usuários e ferramentas de inteligência artificial, tornando as respostas mais úteis e relevantes. Além disso, contribui para o uso responsável dessas tecnologias, ajudando a prevenir a geração de conteúdos inadequados.

Entre as técnicas utilizadas estão a cadeia de pensamento, que consiste em estruturar instruções para orientar a IA na resolução de um problema, e a engenharia de prompt reversa, na qual o usuário apresenta um resultado desejado e solicita à ferramenta que elabore um prompt capaz de produzi-lo. Essas técnicas ajudam a explorar melhor os recursos da inteligência artificial e a obter resultados mais próximos do esperado.

\subsection*{Prompts robustos e testados}
Com tudo o que disse não espere encontrar prompts incompletos como: \textit{"Faça uma auditoria completa do meu CV final: [cole]. Verifique cada seção: compatibilidade com ATS, força das palavras-chave, qualidade das conquistas, clareza da formatação, impacto do resumo, relevância das habilidades, adequação do tamanho. Dê uma nota de 0 a 10 para cada seção. Diga-me exatamente o que ainda está fraco e a correção precisa para cada item. Quero que este CV tire 90+ antes de enviá-lo para qualquer lugar"}.

É ótimo para mostrar conhecimento em redes sociais e conseguir \textit{likes} de ingênuos, mas é um prompt péssimo, existem uma série de erros nele, sete erros para ser mais específico, sendo um deles grave:
\begin{enumerate}
	\item A "nota 90+" volta pela porta dos fundos. Não existe score único de ATS, mas aqui a meta aparece como "quero que este CV tire 90+". O modelo vai inventar notas por seção para satisfazer o pedido, e vai otimizar para uma métrica fictícia. Pior: notas numéricas criam falsa precisão - um "7/10 em palavras-chave" parece medido, mas é atribuído arbitrariamente.
	\item Sete dimensões auditadas ao mesmo tempo, sem hierarquia. Compatibilidade ATS, palavras-chave, conquistas, formatação, resumo, habilidades e tamanho têm pesos diferentes e ordens de prioridade. Se o formato quebra o parsing de ATS, pouco importa que as conquistas sejam boas. Sem hierarquia, o modelo distribui atenção igualmente e não sabe o que corrigir primeiro.
	\item Falta o cargo-alvo. Auditoria de CV sem vaga-alvo é auditoria no vácuo. "Impacto do resumo" e "relevância das habilidades" só fazem sentido contra um alvo. Sem isso, o modelo audita o CV como objeto genérico — e o veredito é inútil.
	\item Notas numéricas por seção não são acionáveis. Saber que "formatação = 7/10" não diz o que mudar. A correção precisa ser o defeito específico + a reescrita, não a nota. A nota é o resumo do diagnóstico, não o diagnóstico.
	\item "Faça uma auditoria completa" sem checklist. "Completa" é vago. O que compõe uma auditoria? Sem checklist explícito, cada execução audita coisas diferentes - e não se consegue comparar versões do CV ao longo do tempo.
	\item Risco de fabricação no "correção precisa". Se o modelo sugere "adicione um número ao bullet 3", ele pode sugerir um número - o mesmo problema fatal do prompt D. A trava de [CONFIRMAR] precisa estar aqui.
	\item Falta a pergunta mais importante de uma auditoria final: "o CV está pronto para enviar?" O prompt termina em "o que ainda está fraco", sem um veredito de prontidão. Você sai com uma lista de problemas e sem saber se pode enviar hoje ou se precisa de outra rodada.
	\item Há ainda um problema silencioso: o modelo não conhece o histórico do CV. Sem saber em que versão você está (primeira? décima?), nem o que já tentou, a auditoria trata tudo como se fosse a primeira vez - e você pode receber sugestões que já aplicou.
\end{enumerate}

Teremos prompts completos, bem escritos e principalmente, testados. Não se preocupe com o tamanho de cada prompt ou qual IA deve usar (qualquer uma servirá e entenderá muito bem o pedido).

\section*{Turbinar minhas leituras}
\addcontentsline{toc}{section}{Prompts}
\begin{promptbox}[A. Entender o livro antes de começar]
Aja como um crítico literário e professor de literatura. Vou ler [título], de [autor], e quero um panorama preparatório antes de começar. Organize a resposta nas seguintes seções, nesta ordem:

O autor em 3-5 linhas - quem é, época, estilo, lugar na tradição literária.

Contexto histórico, social e cultural - o que estava acontecendo quando a obra foi escrita e como isso a molda.

Temas principais e ideias centrais - liste de 3 a 5, cada um com 1-2 frases.

Personagens principais (se for ficção) ou conceitos-chave (se for não-ficção) - com breve descrição de cada.

O que devo ter em mente ao ler - 2-3 dicas de leitura (ex.: estrutura incomum, simbolismo recorrente, dificuldade de linguagem).

Regras: seja direto, sem enrolação; use linguagem acessível; se não tiver certeza de algum fato, sinalize em vez de inventar. Extensão total: até ~500 palavras."
\end{promptbox}

\begin{promptbox}[B. Criar um cronograma de leitura]
Aja como um planejador de estudos. Quero um cronograma de leitura para [título], de [autor].

Dados:

Total de páginas: [x]

Prazo desejado: [x] dias

Minha disponibilidade: [ex.: 40 min por dia, de segunda a sexta] ou [ex.: 1h aos fins de semana]

Nível de dificuldade que espero: [leve / médio / denso]

Instruções:

Divida a leitura em blocos diários, indicando páginas (e capítulos, se você souber a divisão da edição; se não souber, diga que está estimando por páginas).

Se o livro tiver capítulos de tamanhos muito diferentes, não corte capítulos no meio - ajuste o número de páginas do dia.

Inclua 1 dia de folga a cada 6 dias de leitura (ou 2, se o prazo permitir) para não acumular atraso.

Apresente como tabela: Dia | Páginas | Capítulos | Observação (ex.: 'capítulo denso', 'trecho descritivo').

No fim, dê 2 dicas para manter o ritmo (ex.: o que fazer se atrasar, como marcar o progresso).

Regras: seja realista; se o prazo for apertado demais para o número de páginas, avise antes de montar e sugira um prazo alternativo. Não invente números de capítulos se não tiver certeza.
\end{promptbox}

\begin{promptbox}[C. Esclarecer dúvidas durante a leitura]
Aja como um professor de literatura explicando um trecho para um leitor atento, mas não especialista. Analise o trecho abaixo de [título], de [autor]:

'[trecho]'

Contexto que já tenho: [ex.: já li até o capítulo X / é a primeira vez que leio o autor / conheço o enredo geral].

Estruture a resposta assim:

Paráfrase - o que está sendo dito, literalmente, em 2-3 frases.

O que o texto sugere - evite afirmar 'a intenção do autor'; fale do que o trecho permite ler (ambiguidades, escolhas de linguagem, tom).

Duas ou três interpretações possíveis - cada uma com 1-2 frases e o que no trecho a sustenta.

Conexões - no máximo duas: uma com outra obra literária e uma com um conceito teórico/filosófico. Só inclua se a conexão for realmente pertinente; se não houver, diga isso.

Conexão com o tema central - como este trecho ecoa ou complica o tema maior da obra.

Regras: ancore cada afirmação em uma palavra, imagem ou estrutura do próprio trecho; não invente citações; se o trecho for ambíguo, diga que é ambíguo em vez de resolver a ambiguidade. Linguagem clara, sem jargão desnecessário. Extensão: até ~400 palavras.
\end{promptbox}

\begin{promptbox}[D. Organizar suas anotações]
Aja como um assistente de estudos literários. Vou colar abaixo minhas anotações e citações do livro [titulo], de [autor]. Organize-as agrupando por tipo e intenção - por exemplo: falas de personagens, descrições de cenário, diálogos marcantes, passagens temáticas, trechos que geraram dúvida. Para cada grupo: (1) apresente as citações em lista; (2) acrescente uma linha explicando por que elas se agrupam ali; (3) indique a página ou capítulo, se eu tiver fornecido. Se alguma citação puder pertencer a mais de um grupo, coloque-a no grupo principal e mencione os secundários entre parênteses. Não reescreva nem "melhore" minhas citações - preserve o texto original. No fim, sugira 2 ou 3 categorias adicionais que eu não usei, mas que poderiam ser úteis. Formato: títulos de seção + listas. Se eu não tiver fornecido página/capítulo, apenas omita essa informação, sem inventar.
\end{promptbox}

\begin{promptbox}[E. Preparar-se para discutir ou escrever sobre o livro]
Aja como um professor que vai me arguir sobre [titulo], de [autor]. Com base no que li (e nas anotações que eu colar, se eu colar), monte um roteiro de preparação com: (1) 5 perguntas de compreensão que uma discussão séria exigiria, com uma resposta-síntese de 2-3 frases cada; (2) 3 perguntas abertas ou polêmicas, sem resposta pronta, para eu pensar; (3) 3 armadilhas comuns de leitura deste livro (interpretações fáceis ou enganosas) e por que são armadilhas; (4) uma lista de 5 passagens-chave que eu deveria saber localizar, com indicação do que cada uma demonstra. Regras: não invente citações literais; se não tiver certeza de uma passagem, descreva-a em vez de citá-la. Se eu não colar anotações, baseie-se no conhecimento geral da obra e sinalize isso no início. Linguagem direta, sem jargão. Extensão: até ~600 palavras.
\end{promptbox}

\begin{promptbox}[F. Produzir resumos e resenhas]
Aja como crítico literário escrevendo uma resenha para publicação. Com base nas minhas anotações e marcações abaixo sobre [titulo], de [autor]] produza uma resenha com esta estrutura:

Abertura (2-3 frases) - uma cena, imagem ou questão que fisgue o leitor, sem entregar o enredo.

O livro em poucas linhas - do que trata, onde se passa, quem protagoniza. Regra anti-spoiler: não revele reviravoltas, mortes, revelações de identidade ou o desfecho; limite-se ao que aparece no primeiro terço ou ao que a contracapa já informa.

Personagens e relações - análise de 2 a 4 personagens centrais, com foco no que eles representam e como se relacionam, não no que fazem.

Temas centrais - 3 temas, cada um conectado explicitamente a um trecho que eu marquei (cite o trecho entre aspas, se eu o tiver fornecido; caso contrário, descreva-o).

Avaliação da escrita - o que a prosa do autor faz bem, com 1 exemplo concreto (ritmo, diálogo, descrição, estrutura).

Minha impressão final - NÃO invente. Em vez de escrever a minha impressão, deixe um espaço marcado como: [SUA IMPRESSÃO AQUI: escreva 3-5 frases sobre o que sentiu ao terminar o livro]. Eu preencho depois.

Regras: tom envolvente, mas sem adjetivos vazios ("incrível", "emocionante") sem justificativa; nada de spoilers; se eu não fornecer anotações suficientes para alguma seção, sinalize em vez de preencher com generalidades. Extensão: ~500-700 palavras.
\end{promptbox}

\subsection*{Dicas úteis nessa seção}
Nunca peça ao modelo um conteúdo que só existe na sua cabeça (sua impressão, sua opinião, sua experiência). Peça que ele deixe o espaço preparado para você preencher, ou que faça perguntas para extrair isso de você.

Regra geral, que vale para todo prompt de \textit{"me dê ideias de X"}. Não peça uma lista. Peça uma sequência com relação explícita entre os itens. Listas são fáceis de gerar e fáceis de esquecer. Sequências obrigam o modelo a pensar em progressão - e obrigam você a executar em ordem.

\section*{Plano de Estudo}
\addcontentsline{toc}{section}{Prompts}

\begin{promptbox}[A. Criar um plano de aprendizado personalizado]
Aja como um designer instrucional experiente em aprendizagem de adultos. 
Quero um plano de aprendizado para [habilidade], do nível iniciante absoluto até um objetivo definido.

Meus dados:
  Conhecimento atual: [descreva o que você já sabe]
  Tempo disponível: [horas por semana], distribuídas em [ex.: 3 sessões de 1h]
  Estilo de aprendizado: [visual/prático/leitura/auditivo]
  Objetivo final, definido de forma observável: [ex.: "conversar 15 min em espanhol sobre temas do dia a dia" / "construir um app CRUD sozinho"]
  Prazo desejado: [x semanas]

Antes de montar o plano:
  Avalie se o objetivo é realista no prazo com o tempo disponível. Se não for, diga isso primeiro e proponha duas alternativas: (a) manter o objetivo e estender o prazo; (b) manter o prazo e reduzir o escopo do objetivo. Só monte o plano depois que eu escolher.
  Se você não tiver conhecimento confiável dos recursos específicos dessa habilidade, diga isso e trabalhe com princípios gerais de aprendizagem em vez de inventar nomes de cursos, livros ou ferramentas.

Estrutura do plano (por semana):
  Foco da semana (1 frase)
  Tópicos em ordem, com tempo estimado para cada um
  Atividade principal, adaptada ao meu estilo de aprendizado ([visual/prático/leitura]) - explique como ela reflete esse estilo
  Marco verificável - algo que eu consiga fazer ou mostrar ao fim da semana, não apenas "entender"
  Métrica de progresso - como sei que avancei (ex.: gravar 1 min de fala, resolver 5 exercícios, escrever 1 resumo)

Regras adicionais:
 A cada 4 semanas, inclua uma semana de consolidação e revisão, não de conteúdo novo.
 Inclua um plano de contingência: o que fazer se eu atrasar 1 semana (o que cortar primeiro, o que nunca cortar).
 No fim, liste 3 sinais de que o plano precisa ser refeito (ex.: objetivo mudou, tempo caiu pela metade, dificuldade acima do esperado).

Formato: tabela semanal + seções curtas de orientação. Extensão: o mais enxuto possível sem perder detalhe.
\end{promptbox}

\begin{promptbox}[B. Ter conceitos explicados de várias maneiras]
Aja como um professor que domina tanto o conceito quanto a arte de explicá-lo. Quero entender [conceito].

Meu contexto:
 Por que quero entender isso: [ex.: preciso aplicar no trabalho / curiosidade / pré-requisito para outro assunto]
 Onde eu travei: [ex.: não entendo o que significa X / entendo a definição mas não vejo utilidade / confundo com Y]
 O que já sei sobre o tema: [breve]

Explique em 4 camadas, nesta ordem, cada uma construindo sobre a anterior:
 A intuição - a ideia central em linguagem cotidiana, sem jargão, em 2-3 frases. (Não "como se eu tivesse 10 anos" - como se eu tivesse pressa e nenhum vocabulário técnico.)
 A analogia - uma comparação com algo do mundo real que capture a estrutura do conceito (não apenas a superfície). Diga explicitamente onde a analogia funciona e onde ela quebra.
 A explicação técnica - a definição precisa, com os termos corretos, as partes que compõem o conceito e como elas se relacionam. Aqui pode usar jargão, mas defina cada termo na primeira aparição.
 O cenário concreto - uma situação específica (real ou hipotética) em que o conceito aparece em ação, com detalhes suficientes para eu reconhecê-lo quando encontrar algo parecido.

Regras:
 As quatro camadas devem ser distintas, não quatro reformulações da mesma ideia.
 Não repita entre camadas; cada uma acrescenta algo que a anterior não tinha.
 Se o conceito tiver variações ou escolas diferentes (ex.: definição A vs. definição B), mencione isso na camada 3.
 Se você não tiver certeza sobre algum ponto, sinalize em vez de preencher com plausibilidade.

Ao final, para verificar se eu entendi:
 Faça 2 perguntas que eu deveria conseguir responder se a explicação funcionou.
 Indique 1 erro comum que pessoas cometem ao aprender esse conceito.
 Formato: as 4 camadas como seções numeradas, seguidas das duas perguntas e do erro comum. Extensão: até ~500 palavras.
\end{promptbox}

\begin{promptbox}[C. Criar projetos práticos que importam]
Aja como um mentor que já orientou dezenas de pessoas a construir portfólios reais na área de [habilidade].

Meus dados:
 Habilidade: [habilidade]
 Nível atual: [iniciante/intermediário/avançado], com [breve descrição do que já consigo fazer sozinho]
 Interesses: [o que me motiva - ex.: jogos, música, saúde, finanças pessoais, educação]
 Contexto real onde posso aplicar: [ex.: meu trabalho, minha comunidade, um problema que enfrento]
 Tempo disponível por projeto: [horas], distribuídas em [ex.: 2 sessões por semana]

Crie 5 ideias de projetos com progressão real de dificuldade. Cada projeto deve declarar explicitamente:
 Título e problema que resolve - em 1 frase, algo concreto (não "praticar X", mas "resolver Y").
 Por que é adequado ao meu nível - o que ele exige que eu já saiba e o que ele introduz de novo.
 O que muda em relação ao projeto anterior - a progressão não é só "mais difícil"; é: nova técnica? novo escopo? nova restrição (prazo, qualidade, autonomia)? Seja explícito.
 Entregável visível - o que eu mostro a alguém em 2 minutos para provar que fiz. Se não houver entregável claro, o projeto não é de portfólio.
 Etapas em ordem, com estimativa de tempo realista (conte pesquisa, erros e finalização; quando em dúvida, superestime).
 O que vou aprender - habilidade técnica + habilidade prática (ex.: tomar decisão de escopo, documentar, apresentar).
 Definição de "pronto" - critério objetivo de conclusão. Ex.: "roda sem erro em 3 cenários", "tem README", "gravei demo de 1 min".

Regras adicionais:
 Os 5 projetos devem acumular: cada um reaproveita algo do anterior (código, aprendizado, material) ou aprofunda um aspecto dele. Não quero 5 projetos soltos.
 Se algum projeto exigir um recurso que eu não informei ter (ferramenta paga, hardware, acesso a dados), sinalize e ofereça alternativa.
 Inclua 1 projeto "teto": o mais ambicioso que ainda cabe no meu tempo, para eu saber onde está o limite.
 Ao final, diga qual dos 5 eu deveria começar primeiro e por quê.

Formato: um bloco por projeto, com os 7 itens nomeados. Extensão: o mais enxuto possível sem perder o detalhe acionável.	
\end{promptbox}

\begin{promptbox}[D. Obter respostas instantâneas para perguntas específicas]
Aja como um tutor que responde dúvidas pontuais com precisão cirúrgica. Estou travado em [habilidade/tópico].

Onde estou:
 O que eu não entendo: [descreva exatamente o ponto]
 Como eu descreveria meu travamento: [ ] não entendo o que isso significa / [ ] entendo o que é, mas não sei quando ou por que usar / [ ] sei usar, mas o resultado sai errado ou inconsistente / [ ] entendo isolado, mas não vejo como se conecta ao resto
 O que já tentei: [ex.: li a definição, vi um vídeo, tentei aplicar em X e deu Y]

Antes de responder:
 Reformule minha pergunta em uma frase, do jeito que você a entendeu. Se a minha pergunta estiver mal formulada ou partir de uma premissa errada, diga isso primeiro e proponha a pergunta correta.

Se faltar informação para responder com precisão, faça até 2 perguntas antes de prosseguir. Se não precisar, siga direto.

Depois, responda nesta ordem:
 A ideia central - em 2-3 frases, o que resolve meu travamento. Sem rodeios.
 Por que funciona assim - o mecanismo ou a lógica por trás. Aqui, conecte explicitamente ao tipo de travamento que eu marquei acima.
 Erros comuns no meu ponto exato - não os erros genéricos do tema, mas os erros de quem está travado onde eu estou. Se houver mais de um, ordene do mais frequente ao menos.
 Um exemplo isomórfico - mesma estrutura do meu caso, mas com contexto, números ou nomes diferentes, para eu poder consultar sem copiar. Depois, aponte o que no exemplo corresponde ao quê no meu caso.

Para eu verificar se entendi:
 Faça 1 pergunta de aplicação que eu deveria conseguir responder se a explicação funcionou (não uma pergunta de memória; uma que exija usar o conceito).
 Indique 1 sinal de que ainda não entendi - um sintoma que eu deveria observar em mim mesmo.

Regras: se você não tiver certeza sobre algum ponto, sinalize em vez de preencher com plausibilidade. Não repita a mesma ideia em itens diferentes. Formato: seções numeradas como acima. Extensão: até ~400 palavras.
\end{promptbox}

\begin{promptbox}[E. Criar flashcards e materiais de estudo personalizados]
Aja como um designer de materiais de estudo que trabalha com repetição espaçada e prática de recuperação. Quero um guia de estudos para [tópico].

Meus dados:
 Objetivo: [prova / certificação / uso prático / entender de verdade]
 Contexto de uso: [ex.: vou revisar em sessões de 20 min / preciso consultar no trabalho / é pré-requisito para outro assunto]
 O que já domino do tópico: [breve]
 Onde costumo errar: [ex.: confundo termos parecidos / esqueço fórmulas / entendo mas não aplico]

Estrutura do guia:
 Parte 1 - Mapa de conceitos (não fixe em 20; deixe o tópico determinar)
  Liste os conceitos realmente memorizáveis do tópico, em número adequado (entre 10 e 25, se o tópico permitir). Para cada um: nome curto + 1 frase de definição. Ordene por dependência (o que precisa vir antes de quê). Marque com * os 5 mais críticos.

 Parte 2 - Flashcards (apenas para os conceitos do mapa)
  Para cada conceito, 1 ou 2 flashcards no formato:
   Frente: pergunta curta que testa uma única coisa
   Verso: resposta em até 2 linhas, sem parágrafos
  Tipo: definição / aplicação / distinção / exceção (varie os tipos, não faça 20 de definição)
  Regras: nada de perguntas com "explique" ou "discuta"; nada de respostas longas. Se um conceito não render um bom flashcard, diga isso em vez de forçar.

 Parte 3 - Mnemônicos
  Só para os conceitos marcados com *. Para cada: o mnemônico + por que ele funciona (que associação ele explora). Se não houver mnemônico natural para um conceito, diga isso em vez de inventar um forçado.

 Parte 4 - Equívocos ancorados
  Liste os equívocos mais comuns, cada um atrelado ao conceito específico onde ele ocorre (ex.: "Conceito 7: é comum achar que X, mas na verdade Y"). Ordene do mais frequente ao menos.

 Parte 5 - Autoavaliação (diferente dos flashcards)
  5 a 8 perguntas que testem os mesmos conceitos, mas em contexto de aplicação - situações, casos, comparações, não definições repetidas. Cada pergunta com resposta comentada (não só a resposta; por que aquela resposta).

Regras adicionais:
 Se o tópico não tiver 20 conceitos memorizáveis, não infle - diga quantos tem de fato.
 Se você não tiver certeza sobre um mnemônico, equívoco ou conceito, sinalize em vez de preencher.

Formato: seções numeradas, com os conceitos numerados de forma consistente em todas as partes (o conceito 7 na Parte 1 é o mesmo na Parte 4).

Extensão: enxuta o suficiente para ser usada em sessões curtas. Prefira corte a verbosidade.

Ao final:
 Diga em que ordem estudar (quais conceitos primeiro, quais revisar juntos).
 Indique 1 sinal de que eu só decorei, mas não entendi - um sintoma observável.
\end{promptbox}

\begin{promptbox}[F. Aprender por meio da técnica de Feynman]
Aja como um tutor que aplica a Técnica de Feynman de verdade - em ciclo, não em uma resposta só. Quero entender [conceito].

Contexto:
 Por que quero entender: [ex.: prova / aplicar no trabalho / pré-requisito]
 O que já sei: [breve]
 Onde suspeito que vou travar: [ex.: confundo com X / não entendo o porquê]
 Como vamos proceder - em etapas separadas, uma por vez:

Etapa 1 - Você me dá o alvo, não a explicação.
 Antes de explicar qualquer coisa, liste os elementos que uma explicação completa deste conceito deveria conter (ex.: definição, mecanismo, condições de uso, exceções, relação com conceitos vizinhos). Isso vira o critério contra o qual você vai avaliar minha explicação depois. Não explique o conceito ainda.

Etapa 2 - Eu explico primeiro.
 Você me pede para explicar [conceito] com minhas próprias palavras, como se estivesse ensinando alguém que não sabe nada. Me dá um formato mínimo para eu seguir (ex.: "diga o que é, para que serve, como funciona, e dê um exemplo"). E para aqui - espere minha resposta.

Etapa 3 - Você diagnostica, usando o alvo da Etapa 1.
 Quando eu enviar minha explicação, avalie contra os elementos listados na Etapa 1. Para cada lacuna, classifique:
  Ausente - não mencionei
  Errado - mencionei, mas incorretamente
  Vago - mencionei, mas sem precisão
  Desconectado - certo, mas sem ligação com o resto
 Não elogie por educação; seja específico. Se eu acertei algo, diga o que foi e por que está certo.

Etapa 4 - Você reensina apenas o que faltou.
 Não repita a explicação inteira. Reensine só as lacunas da Etapa 3, na ordem em que aparecem. Para cada uma: explique de forma simples, dê um exemplo, e diga como saber que entendi (um sinal observável).

Etapa 5 - Fechamos o ciclo.
 Peça que eu explique de novo apenas os pontos que você reensinou. Se eu ainda errar, repita a Etapa 4 só nesses pontos. Só encerre quando eu conseguir explicar sem os erros anteriores.

Regras:
 Nunca invente a minha explicação. Se eu não respondi ainda, pare e espere.
 Se você não tiver certeza sobre algum ponto do conceito, sinalize em vez de preencher.
 Se o conceito tiver mais de uma interpretação ou escola, mencione isso na Etapa 1 (entra no alvo).
 Linguagem simples nas Etapas 1 e 4; precisão técnica nas Etapas 3 e 5.

Formato: uma etapa por resposta, sem adiantar as seguintes.

Comece agora pela Etapa 1.
\end{promptbox}

\subsection*{Dicas úteis nessa seção}
O maior risco deste prompt não é o modelo errar - é ele acertar demais. Quatro explicações bonitas e coerentes podem produzir a ilusão de entendimento sem que a ficha caia de fato. Isso se chama fluência ilusória: você lê, acha claro, e não consegue aplicar.

Por isso a correção mais importante foi adicionar verificação ao final (as 2 perguntas e o erro comum). Sem isso, o prompt é confortável, mas não necessariamente eficaz.

A mesma lógica vale para o conjunto: todo prompt de aprendizado deveria terminar com um teste de compreensão, não apenas com a entrega do conteúdo. No prompt A isso aparece como "marco verificável"; aqui, como perguntas de checagem.

\section*{Correção do meu CV}
\addcontentsline{toc}{section}{Prompts}

\begin{promptbox}[A. Analisador de raio-X do CV]
Aja como um recrutador sênior que já triou milhares de CVs para vagas de [cargo/área] em [indústria/setor]. Vou anexar meu CV abaixo e quero uma análise operacional, não motivacional.
	
Contexto obrigatório (se eu não preencher, pergunte antes de analisar):
 Cargo-alvo: (ex.: analista de dados pleno)
 Descrição da vaga (se houver): [cole os requisitos principais]
 Minha situação: [ex.: empregado buscando transição / desempregado / recém-formado]
 Setor: [ex.: tecnologia, saúde, finanças - ou "transição entre setores"]
	
Antes de analisar:
 Se algum campo de contexto estiver em branco, pergunte antes de prosseguir.
 Não presuma cargo nem setor.
 Se o CV não tiver informação suficiente para avaliar algum eixo, diga isso em vez de preencher com suposições.

Analise em quatro eixos separados:
 Eixo 1 - Triagem de 10 segundos (o que salta ou afunda)
  Liste os sinais concretos que um recrutador capta na primeira passada: cargo atual legível? anos de experiência visíveis? cargo-alvo reconhecível? formatação legível? Se houver risco de descarte, diga qual sinal o dispara e por quê - sem falar em "seções", fale em elementos específicos.

 Eixo 2 - Forma (layout, ordem, extensão)
  Avalie: extensão adequada ao nível? ordem das seções condiz com o esperado para o cargo-alvo? legibilidade (fonte, espaçamento, bullets)? consistência (datas, tempos verbais, formatação)? Aponte o que muda e o que fica.
	
 Eixo 3 - Conteúdo (conquistas e relevância)
  Para cada experiência listada no CV, avalie:
   A conquista é resultado ou tarefa? (tarefa = "responsável por X"; resultado = "reduzi X em Y%")
   Tem número, escala ou consequência?
   É relevante para o cargo-alvo ou é ruído?
   Classifique cada bullet como: forte / fraca / irrelevante / ambígua. Para as fracas e ambíguas, reescreva em uma versão melhor (não diga "melhore"; mostre como).
	
 Eixo 4 - Riscos de rejeição (não "a única coisa")
  Liste todos os riscos plausíveis de rejeição, ordenados por gravidade. Para cada um: qual é o risco, por que um recrutador se incomodaria, e como mitigar (o que reescrever, mover, ou remover). Se houver apenas um risco relevante, diga que é apenas um - não force um único motivo se houver vários, e não invente um se não houver.
	
Priorização (o mais importante):
 Ao final, monte uma lista das 5 correções com maior retorno, ordenadas por impacto. Para cada uma: o que mudar, onde no CV, e o efeito esperado na triagem.
	
Regras:
 Seja específico: aponte o bullet, a linha, a seção exata. Crítica genérica não serve.
 Não use adjetivos de efeito ("desastroso", "incrível") sem justificar com o critério.
 Se algo estiver bom, diga o que está bom e por quê - para eu saber o que não mexer.
 Se você não tiver certeza sobre convenções do setor/região, sinalize em vez de afirmar.
 Formato: quatro eixos + priorização. Sem introdução motivacional.
\end{promptbox}

\textbf{ATENÇÃO}: Não esqueça de adicionar seu currículo.

\begin{promptbox}[B. Decodificador da descrição da vaga]
Aja como um recrutador experiente que sabe ler vagas nas entrelinhas. Vou colar a descrição de uma vaga e quero uma decodificação ancorada no texto, não em suposições sobre a empresa.
	
Insumos obrigatórios (se eu não preencher algum, pergunte antes de analisar):
 Descrição da vaga: [cole a descrição aqui]
 Empresa e setor: [nome, se souber / setor / tamanho aproximado]
 Meu perfil resumido (opcional, mas recomendado): [ex.: 5 anos em X, mudei de Y para Z, tenho experiência em A e B]
 Meu objetivo: [ex.: passar na triagem / me preparar para entrevista / decidir se vale a pena me candidatar]
	
Antes de decodificar:
 Se a descrição for curta ou genérica demais para inferências confiáveis, diga isso primeiro e trabalhe só com o que o texto sustenta.
 Se você não tiver informação sobre a empresa, não invente cultura, valores ou prioridades; marque essas inferências como "genéricas do setor", não da empresa.
	
Decodifique em cinco camadas:

1. O que a vaga diz explicitamente
 Separe em: requisitos obrigatórios (must-have) e desejáveis (nice-to-have). Cite o trecho exato de cada um. Não parafraseie - transcreva.
	
2. O que a vaga implica (habilidades ocultas)
 Para cada inferência, mostre o trecho que a sustenta e o tipo de raciocínio:
  Repetição (ex.: "autonomia" aparece 3x -> espera-se alguém que não precise de supervisão)
  Verbo escolhido (ex.: "estruturar" vs. "manter" indica nível de senioridade)
  Omissão (ex.: não menciona equipe -> possível trabalho solo ou liderança não dita)
  Contexto do setor (ex.: "ritmo acelerado" em startup significa algo diferente de em banco)
  Classifique cada inferência por confiança: alta / média / baixa.
	
3. Sinais de alerta para um candidato
 Liste os pontos em que um candidato sem determinada experiência ou sinal tende a ser descartado - sempre ancorado em um requisito ou inferência específica da camada anterior. Se não houver sinal claro, diga que não há em vez de inventar.
	
4. Palavras-chave para o CV
 Não fixe em 3. Liste as palavras/expressões que precisam aparecer no CV para passar na triagem automática e humana. Para cada uma:
  O termo exato (como aparece na vaga)
  Onde deve aparecer (resumo, experiência, habilidades)
  Se é obrigatória ou reforço
  Como você apareceria se eu já tenho essa experiência (dê uma formulação de exemplo)
  Se eu fornecer meu perfil, cruze: quais dessas palavras eu já demonstro e quais preciso evidenciar melhor.
	
5. O que o gerente realmente valoriza (e a vaga não diz)
 Aqui separe duas coisas:
  Inferível do texto: o que a vaga sugere sobre prioridades (ex.: ênfase em "resultados mensuráveis" -> espera-se alguém orientado a métricas).
  Não inferível: o que só se sabe com informação sobre a empresa, o time ou a cultura. Se não tiver essa informação, diga que não é possível saber e indique como descobrir (ex.: pesquisar no LinkedIn, perguntar na entrevista, olhar reviews).
	
Ao final:
 3 a 5 ajustes concretos no meu CV/material para esta vaga, em ordem de impacto.
 2 ou 3 perguntas que eu deveria fazer na entrevista para confirmar as inferências de confiança baixa.
 1 sinal de que esta vaga pode não valer a pena - um indício no texto de que o papel ou a empresa têm um problema (ex.: acúmulo de funções, requisitos contraditórios, linguagem que revela alta rotatividade).
	
Regras:
 Toda inferência precisa citar o trecho que a origina. Sem trecho, é suposição - e deve ser marcada como tal.
 Não use generalidades de vaga ("querem alguém proativo") sem ancoragem no texto.
 Se algo não for decodificável, diga "não é possível inferir" em vez de preencher.
 Formato: cinco camadas numeradas + fechamento. Sem introdução motivacional.
\end{promptbox}

\begin{promptbox}[C. Otimizador para bypass de ATS]
Aja como alguém que entende como ATS funcionam na prática - incluindo seus limites. Não finja ter acesso a um ATS real nem atribua notas numéricas; trabalhe com princípios verificáveis de correspondência entre CV e vaga.
	
Insumos obrigatórios (se eu não preencher, pergunte antes):
 Meu CV atual: [cole o CV]
 Descrição da vaga: [cole a vaga]
 País/região da candidatura: [ex.: Brasil, EUA, Europa] - as convenções de ATS variam
 Plataforma, se souber: [ex.: Gupy, Workday, LinkedIn Easy Apply]

Antes de otimizar:
 Deixe claro que não existe nota única de ATS e que você não pode medir minha pontuação. Se o meu CV é reprovado, as causas prováveis são: formato ilegível, ausência de termos-chave, falta de correspondência com a vaga, ou problemas de estrutura - analise por aí, não por score inventado.
	
Se você não tiver certeza sobre a convenção de ATS da minha região/plataforma, sinalize em vez de afirmar.
	
Etapa 1 - Diagnóstico de correspondência (não de score)
Liste:
 Termos da vaga que aparecem no meu CV - com onde aparecem
 Termos da vaga que não aparecem - e classifique cada um como:
 Tenho experiência, falta o termo (corrigível com reescrita)
 Tenho algo próximo, mas não equivalente (corrigível com reformulação honesta)
 Não tenho (não se resolve com palavra-chave - ver Etapa 4)
 Termos boilerplate (presentes em quase toda vaga, baixo poder de diferenciação) vs. termos diferenciadores (específicos da vaga, alto poder)
	
Etapa 2 - Diagnóstico de formato e estrutura (o que quebra ATS)
 Aponte problemas que realmente atrapalham parsing:
 Colunas, tabelas, caixas de texto, cabeçalhos/rodapés
 Fontes não-padrão, ícones no lugar de texto
 Títulos de seção não convencionais ("Minha jornada" em vez de "Experiência")
 Datas em formato inconsistente
 PDF gerado de imagem (não-parseável)
 Diga o que mudar e por quê. Se o formato estiver OK, diga que está OK.
	
Etapa 3 - Reescrita de Experiência e Habilidades
Para cada alteração:
 Antes (transcreva o bullet original)
 Depois (reescrito)
 O que mudou (termo adicionado? estrutura? verbo? número?)
 De onde vem a informação (a mudança se baseia no que já estava no CV? ou em algo que eu preciso confirmar?)
	
Regras da reescrita:
 Use os termos da vaga onde eles são verdadeiros para mim. Se eu não tenho a experiência, não invente.
 Não empilhe palavras-chave. 
 Não transforme bullets em listas de termos. Um termo por frase, no contexto certo.
 Não invente números, cargos, ferramentas ou responsabilidades. Se um número ou termo precisa vir de mim, deixe [CONFIRMAR: ...] no lugar.
 Preserve o que já está bom. Não reescreva por reescrever.
 Se a vaga pede algo que eu não tenho, diga claramente: isso é lacuna real e precisa ser tratada na Etapa 4.
	
Etapa 4 - Lacunas reais (o que palavra-chave não resolve)
Liste o que a vaga pede e eu não tenho. Para cada item:
 É eliminatória (sem ela, não passo) ou desejável?
 O que posso fazer: buscar formação rápida, destacar algo adjacente, ou aceitar que a vaga não é para mim.
 Se for eliminatória e eu não tiver como cobrir, diga com clareza: candidatar-se provavelmente não vale o esforço - não empurre uma candidatura inviável.
	
Etapa 5 - Verificação final
 O CV reescrito continua legível para humanos (não virou lista de termos)?
 Os termos-chave aparecem naturalmente em contexto?
 Alguma afirmação precisa da minha confirmação antes de eu enviar? Liste tudo com [CONFIRMAR].
	
Regras gerais:
 Não prometa pontuação de ATS. Não invente score. Não simule um ATS que você não tem.
 Se algo não for verificável, diga "não é possível medir".
 Priorize mudanças por impacto na triagem real: formato > correspondência de termos ] força dos bullets.
 Formato: cinco etapas, com antes/depois na Etapa 3. Sem introdução motivacional.
\end{promptbox}

\begin{promptbox}[D. Transformar as conquistas]
Aja como um revisor de CV que sabe transformar tarefas em resultados sem inventar nada. Vou colar minhas descrições atuais e quero que você as reescreva com honestidade - marcando o que precisa da minha confirmação em vez de preencher com números plausíveis.
	
Insumos obrigatórios (se eu não preencher, pergunte antes):
 Cargo-alvo: [ex.: analista de dados pleno em fintech]
 Setor: [ex.: tecnologia, saúde, educação]
 Nível: [júnior / pleno / sênior / liderança]

Minhas descrições atuais: [cole os bullets do CV]
 
Para cada bullet, se você tiver: o problema que existia antes, o que você fez, e o que mudou depois. Se não tiver, tudo bem - deixe em branco.
	
Antes de reescrever:
 Se algum bullet não tiver informação suficiente para virar resultado, não invente. Marque [FALTA INFO: o que eu preciso saber para reescrever isto].
 Se eu não fornecer números, não sugira números plausíveis. Em vez disso, pergunte: "que dado você tem para estimar isso?" - e dê exemplos de onde procurar (relatórios do período, e-mails, métricas do time).
	
Etapa 1 - Diagnóstico de cada bullet
Para cada bullet, classifique:
 Tarefa (descreve o que você fazia, não o que mudou)
 Resultado sem número (mudou algo, mas sem métrica)
 Resultado com número (já é forte - só ajustar forma)
 Vago (não dá para saber o que você fez)
 E diga o que falta em cada um: número? contexto? verbo mais preciso? clareza de escopo?
	
Etapa 2 - Reescrita com três formatos possíveis
 Não force a mesma fórmula em todos. Escolha o formato que cabe no bullet:
 Fórmula completa (quando há número e impacto): verbo + o que fiz + resultado mensurável + impacto
 Fórmula de resultado (quando há resultado, mas impacto é indireto ou desconhecido): verbo + o que fiz + resultado
 Fórmula de escopo (quando não há número, mas há complexidade): verbo + o que fiz + escala/contexto que dá dimensão
	
Para cada bullet:
 Antes: transcreva o original
 Depois: a reescrita, no formato que couber
 O que mudou: o que foi ajustado e por quê
 O que falta: se ficou [CONFIRMAR: número X] ou [FALTA INFO: Y], diga exatamente o que eu preciso buscar
	
Regras da reescrita:
 Não invente números, ferramentas, escopo, cargos ou responsabilidades. Se o dado não veio de mim, vai como [CONFIRMAR] ou [FALTA INFO], nunca como fato.
 Não empilhe adjetivos ("inovador", "estratégico", "robusto") sem conteúdo. Adjetivo sem fato é ruído.
 Não use verbos inflados para o meu nível. Se sou júnior, "liderei" só se for verdade. Prefira verbos compatíveis com a senioridade.
 Não transforme todo bullet em fórmula completa. Se a fórmula não cabe, use outra - ou diga que o bullet deveria ser removido do CV.
 Preserve o que já está bom. Se um bullet já é resultado com número, só ajuste forma, não reescreva.
 Se algum bullet for fraco demais e não tiver conserto, diga com clareza: "este bullet não agrega; considere remover ou substituir por [tipo de conquista]."
	
Etapa 3 - Verbos por nível
 Liste os verbos que usei e diga se são compatíveis com meu nível (júnior/pleno/sênior/liderança). Se algum estiver inflado, substitua. Se algum estiver fraco, sugira alternativa.
	
Etapa 4 - Onde procurar os números que faltam
 Para cada [CONFIRMAR] ou [FALTA INFO], dê uma dica concreta de onde buscar o dado:
 Relatórios do período (quantos clientes atendi?)
 E-mails e atas (que problema resolvi? qual foi o prazo?)
 Métricas do time (qual era o volume antes e depois?)
 Comparação temporal (como era antes de eu entrar?)
 Não sugira valores. Sugira fontes.
	
Etapa 5 - Verificação final
 Todos os bullets são verdadeiros (nada inventado)?
 Os [CONFIRMAR] e [FALTA INFO] estão marcados e visíveis?
 O conjunto não soa artificial (todos os bullets com a mesma estrutura)?
 Algum bullet deveria ser removido em vez de reescrito?
	
Regras gerais:
 O objetivo não é "fazer o recrutador parar de rolar" - é fazer o recrutador entender rápido o que eu entreguei e por que isso importa para a vaga. Especificidade e relevância fazem isso; floreio não.
 Se você não tiver certeza sobre a convenção do meu setor/região, sinalize.
 Formato: cinco etapas, com antes/depois na Etapa 2. Sem introdução motivacional.
\end{promptbox}

\begin{promptbox}[E. Leitor de mentes do recrutador]
Aja como um recrutador experiente simulando a triagem real de uma vaga. Deixe claro desde o início: você não tem acesso à mente de recrutadores reais nem ao pool de candidatos; o que você fará é uma simulação baseada em princípios conhecidos de triagem (clareza, relevância, sinais de risco, evidência de resultado). Trate a saída como hipótese fundamentada, não como veredito.
	
Insumos obrigatórios (se eu não preencher, pergunte antes):
 Descrição da vaga: [cole a vaga]
 Meu CV: [cole o CV]
 Meu nível: [júnior / pleno / sênior / liderança]
 Meu objetivo com esta análise: [ex.: entender por que sou descartado / ajustar antes de enviar / decidir se vale me candidatar]
	
Antes de simular:
 Se o CV for longo ou o contexto for insuficiente, diga que a simulação será parcial.
 Se você não tiver informação sobre a empresa ou o setor, não invente cultura; trabalhe apenas com o que a vaga e o CV permitem inferir.
 Não escreva cena nem monólogo interno. Faça diagnóstico ancorado em elementos concretos do CV e da vaga.
	
Etapa 1 - Triagem em 6 segundos (sinais, não veredito binário)
Liste, na ordem em que um recrutador captaria:
 O que salta imediatamente (cargo legível? anos de experiência? correspondência com o título da vaga? senioridade visível?)
 O que gera atrito (formatação, extensão, ordem, jargão, lacunas inexplicadas)
 O que gera dúvida (algo que faz o recrutador hesitar, não descartar)
 Para cada item: transcreva o elemento do CV que o origina. Sem transcrição, não vale.
	
Etapa 2 - Veredito de triagem (com nuance, não binário)
 Em vez de "selecionado ou descartado", classifique em um dos três:
 Avançaria para leitura detalhada - por quê (cite os sinais)
 Ficaria na pilha intermediária - por quê (o que falta para avançar)
 Seria descartado - por qual sinal concreto (cite o elemento)
 Se o veredito depender de informação que você não tem (ex.: se o recrutador valoriza diploma X), diga isso e apresente os dois cenários.
	
Etapa 3 - Primeira coisa que chama atenção (positiva)
 Diga qual elemento do CV capta atenção primeiro e por que ele funciona para esta vaga. Se não houver nada que se destaque, diga com clareza: "nada se destaca nos primeiros 6 segundos" - o que é, por si só, um diagnóstico.

Etapa 4 - Sinais que enfraquecem (não "a única coisa")
Liste todos os sinais que reduzem as chances, ordenados por gravidade. Para cada um:
 O elemento no CV (transcrito)
 Por que um recrutador se incomodaria
 Se é problema real (ex.: falta de correspondência com a vaga) ou viés de leitura (ex.: formato incomum, lacuna que tem explicação)
 Como mitigar: reescrever, mover, remover, ou contextualizar
 Se houver apenas um sinal relevante, diga que é apenas um. Não force um único motivo quando há vários, e não invente um quando não há.
	
Etapa 5 - O que faria o recrutador ligar (com trava de honestidade)
Descreva a qualidade ou evidência que de fato faria um recrutador querer conversar. Depois, separe:
 O que eu já tenho e não estou destacando - reescreva o bullet para evidenciar
 O que eu não tenho - diga com clareza que isso não se resolve com reescrita; resolve-se com experiência, formação ou escolha de vaga diferente
 O que você não pode afirmar - se depender de informação sobre a empresa/recrutador, sinalize
 
Proibição explícita:
 Não sugira que eu invente experiência, números, cargos ou ferramentas para me aproximar do que faria o recrutador ligar.
 Se a única forma de chegar lá for fabricar, diga: "isto não se resolve neste CV; requer experiência real que você não tem".
	
Etapa 6 - Verificação final
 Os sinais apontados estão transcritos do CV ou são impressão sem base?
 O veredito é defensável com os elementos citados?
 Há algo que depende de informação externa e foi marcado como tal?
 Alguma sugestão minha exigiria inventar algo? Se sim, remova.
	
Regras gerais:
 Trate a simulação como hipótese fundamentada, não como verdade. Diga isso no início e no fim.
 Não escreva cena. Não finja ser um recrutador específico. Não cite "200 CVs" como se tivesse acesso ao pool.
 Priorize sinais objetivos (correspondência com a vaga, clareza, evidência de resultado) sobre impressões subjetivas (estilo, formato, instituição).
 Formato: seis etapas numeradas. Sem introdução motivacional, sem monólogo interno.
\end{promptbox}

\begin{promptbox}[F. Auditoria final do CV]
Aja como um revisor sênior que faz a auditoria final de um CV antes do envio. Seu trabalho não é elogiar nem dar notas de efeito: é encontrar o que ainda pode derrubar a candidatura e dizer exatamente como corrigir.
	
Insumos obrigatórios (se eu não preencher, pergunte antes):
 Meu CV final: [cole o CV]
 Cargo-alvo: [ex.: gerente de projetos em ONG]
 Nível: [júnior / pleno / sênior / liderança]
 País/região e idioma do CV: [ex.: Brasil, português]
 Versão: [ex.: 3a. revisão, já apliquei correções de ATS e conquistas]
 O que já tentei e não funcionou: [opcional, evita sugestões repetidas]
	
Declare antes de auditar:
 Você não pode atribuir nota real de ATS nem prever aprovação automática. Não dê nota numérica por seção: dê diagnóstico específico e correção.
 Se algum dado do CV não for suficiente para auditar alguma dimensão, diga em vez de preencher.
	
Antes da auditoria:
 Se faltar algum insumo obrigatório, pergunte antes de prosseguir.
 Não invente números, ferramentas, cargos ou experiências ao sugerir correções. Onde faltar dado meu, use [CONFIRMAR: ...].
 Audite as sete dimensões em ordem de prioridade (esta ordem importa):
	
1. Formatação e parseabilidade (a mais crítica)
 É legível por ATS? (colunas, tabelas, caixas, cabeçalhos, fontes, PDF-imagem)
 Títulos de seção convencionais? Datas consistentes?
 Extensão adequada ao nível?
 Se o formato quebra parsing, isso é a primeira coisa a corrigir - nada abaixo importa antes disso.
	
2. Correspondência com o cargo-alvo
 O cargo-alvo é reconhecível nos primeiros 10 segundos?
 As palavras-chave da vaga aparecem onde importam (título, resumo, experiência)?
 Algum termo é boilerplate em vez de diferenciador?
 Há lacunas de correspondência? Classifique: tenho experiência mas falta o termo / tenho algo próximo / não tenho (não se resolve com reescrita).
	
3. Resumo
 Diz o que sou, para que vaga, e com que evidência?
 Ou é genérico ("profissional dedicado com experiência em...")?
 Reescreva se estiver fraco.
	
4. Conquistas
 Para cada bullet: é tarefa, resultado sem número, resultado com número, ou vago?
 Há número, escala ou consequência?
 É relevante para o cargo-alvo ou ruído?
 Aponte os bullets fracos e reescreva - marcando [CONFIRMAR] onde faltar dado meu.
	
5. Habilidades
 Lista é relevante ou virou depósito?
 Separa técnicas de comportamentais? (depende da convenção do país)
 Há habilidade que aparece na vaga e não está no CV?

6. Consistência interna
 Datas, tempos verbais, formatação de bullets, nomes de empresas/cargos consistentes?
 Alguma contradição entre seções (ex.: resumo diz X, experiência mostra Y)?
	
7. Legibilidade humana
 Um humano entende a trajetória em uma passada?
 Há algo que soa artificial (todos os bullets com a mesma fórmula, adjetivos empilhados, verbos inflados para o nível)?
	
Para cada dimensão, entregue:
 Status: OK / precisa ajuste / crítico
 O que está específico (transcreva o elemento do CV)
 Por que é problema (consequência concreta na triagem)
 Correção exata (o que mudar, onde, como fica depois)
	
Depois da auditoria:
 A. Priorização
  Liste as correções em ordem de impacto na triagem, não em ordem de seção. Para cada uma: o que muda, onde, e o efeito esperado. Se houver mais de cinco correções, agrupe as menores.

 B. Itens [CONFIRMAR]
  Liste tudo o que precisa da minha confirmação antes de eu enviar - números, datas, escopo, ferramentas. Não sugira valores; diga o que eu preciso checar.
 
 C. Veredito de prontidão
  Responda com clareza uma das três:
   Pronto para enviar - sem correções críticas pendentes
   Pronto após ajustes menores - liste quais, em ordem
   Não enviar ainda - há problema crítico (formato quebrado, falta de correspondência, lacuna eliminatória); diga qual e o que fazer
	
 Se o veredito for "não enviar ainda" e a lacuna for real (experiência que falta), diga com clareza: isto não se resolve neste CV; requer experiência que você não tem ou vaga diferente.
	
Regras gerais:
 Sem notas numéricas. Sem "90+". Sem score de ATS inventado.
 Toda crítica precisa transcrever o elemento do CV que a origina.
 Toda correção precisa ser acionável - reescrita, movida, removida, ou [CONFIRMAR].
 Se um item já está bom, diga o que está bom e por quê - para eu não mexer.
 Se você não tiver certeza sobre convenção regional/setorial, sinalize.
 Formato: sete dimensões + priorização + veredito. Sem introdução motivacional.
\end{promptbox}

\subsection*{Dicas úteis nessa seção}
A regra geral deste bloco de CV, todo prompt de candidatura deve separar o que o modelo sabe (do documento fornecido), o que ele infere (com base em critérios explícitos), e o que ele não pode saber (empresa, cultura, gerente). Sem essa separação, inferência vira fato, e fato vira estratégia errada.

Dois ajustes opcionais que valem considerar:
\begin{itemize}
	\item Se está decidindo se candidatar, acrescente: "Ao final, dê uma nota de 0 a 10 para meu encaixe na vaga, com base no meu perfil fornecido, e diga o que pesou mais na nota." Isso transforma a decodificação em ferramenta de decisão, não só de preparação.
	\item Se existe a descrição em inglês e vai se candidatar em português (ou vice-versa), acrescente: "Indique também as palavras-chave na versão no outro idioma, respeitando as convenções de CV de cada país." Candidaturas internacionais falham muito por tradução literal de termos.
\end{itemize}

E o fio condutor do bloco inteiro, que vale como fecho: Todo prompt de candidatura pede ao modelo algo que não tem (a mente do recrutador, a nota do ATS, a cultura da empresa, os números que não foram fornecidos). A correção é sempre: declarar o inacessível, ancorar no documento, transcrever o elemento, travar a invenção com [CONFIRMAR], e terminar em veredito acionável.

\section*{Professor particular de idiomas}
\addcontentsline{toc}{section}{Prompts}
1. Parceiro de Conversação \\
"Estou aprendendo [idioma] no nível [nível]. Tenha uma conversa natural comigo inteiramente em [idioma]. Comece no meu nível e aumente gradualmente a dificuldade. Corrija meus erros imediatamente, explique por que estão errados e, depois, continue a conversa."

2. Professor de Gramática \\
"Estou tendo dificuldade com [tópico de gramática] em [idioma]. Explique de forma simples usando exemplos do dia a dia. Mostre erros comuns e, depois, me dê cinco exercícios. Confira minhas respostas e explique o que eu errei."

3. Construtor de Vocabulário \\
"Quero melhorar meu vocabulário de [idioma] em torno de [tópico]. Dê-me 20 palavras e expressões úteis que os falantes nativos realmente usam. Para cada uma, dê o significado, um exemplo e um erro comum. Depois, faça um quiz comigo."

4. Coach de Pronúncia \\
"Eu falo [seu idioma] e estou aprendendo [idioma-alvo]. Identifique os sons que provavelmente serão difíceis para mim. Explique como produzir cada um, os erros comuns e me dê palavras para praticar."

5. Simulador de Falante Nativo \\
"Aja como um falante nativo de [país]. Tenha uma conversa casual comigo sobre [tópico]. Use expressões naturais e frases do dia a dia. Se eu disser algo que soe pouco natural, explique como um falante nativo diria isso."

6. Coach de Fluência \\
"Quero me tornar fluente em conversação em [idioma] dentro de [prazo]. Atualmente estou no nível [nível] e tenho [X horas] por semana. Crie um plano semanal que abranja fala, compreensão auditiva, gramática, vocabulário, pronúncia e leitura. Dê-me tarefas diárias e teste meu progresso."

\section*{Aprender inglês}
\addcontentsline{toc}{section}{Prompts}
1. Amigo de conversa diária \\
"Você é um falante de inglês nativo e amigável. Vamos conversar por 10 minutos em inglês sobre [tópico]. Por favor, corrija meus erros depois de cada mensagem e sugira como eu poderia falar de maneira mais natural." 

2. Detector de erros de gramática \\
"Crie um quiz de 10 perguntas sobre [tema gramatical]. Depois de cada resposta minha, explique o que errei e reforce a regra com um novo exemplo curto." 

3. Turbo vocabulário com memorização \\
"Ensine-me 15 palavras úteis sobre [tema]. Para cada palavra, inclua: significado simples, uma frase do dia a dia e um truque de memória (imagem, rima ou associação)." 

4. Treinador de pronúncia \\
"Analise esta gravação [inserir áudio ou texto para o modelo ler]. Diga onde minha pronúncia precisa melhorar (sons específicos) e me dê 2 exercícios práticos com exemplos para treinar." 

5. Guia de cultura + idioma real \\
"Explique a expressão idiomática [exemplo] usada nos EUA/Reino Unido com: significado, 2 exemplos de uso e uma situação em que seria estranho usar." 

6. Imersão auditiva inteligente \\
"Crie um áudio de 2 minutos em inglês no nível [A1-C1] sobre [tema] e forneça: transcrição, lista de vocabulário difícil e 3 perguntas de compreensão." 

7. 'Médico do texto' \\
"Revise este parágrafo [colar texto]. Depois: a) Corrija meus erros; b) Explique meus 3 principais problemas; c) Mostre a versão nativa/natural do texto." 

8. Criador automático de flashcards \\
"Transforme essas [X] palavras em flashcards de inglês com: pergunta (frase com espaço para completar) + resposta + tradução simples. Formato compatível com Anki." 

9. Plano de imersão personalizado (30 min/dia) \\
"Monte um plano de 4 semanas para eu estudar inglês 30 min por dia, baseado no meu nível [A1-C1]. Inclua links para: vídeos, podcasts, leitura e conversas. Atualize semanalmente com base no meu progresso." 

10. Checkpoint de progresso \\
"Crie uma checklist semanal para me levar do nível [nível atual] para [nível objetivo] em 90 dias, com metas pequenas, recursos sugeridos e um teste rápido semanal para medir evolução." 

\section*{Organizar a alimentação}
\addcontentsline{toc}{section}{Prompts}
1. Busca por tipos de refeições \\
"Eu quero uma lista com [inserir quantidade] refeições adequadas para a minha dieta de [inserir o tipo de dieta]. Sou [inserir informações sobre seu peso, altura, nível de atividade física e outras informações que podem ser relevantes para a dieta]. As refeições devem ser [inserir nível de dificuldade, tempo disponível para o preparo e outras restrições que você pode ter. Ex.: gluten-free]. Mostre refeições variadas e insira o valor nutricional de cada uma delas". 

2. Criação de cardápio semanal \\
"Crie um cardápio semanal adequado para uma dieta de [inserir objetivo]. Vou inserir informações sobre peso, altura, nível de atividade física e outras informações que podem ser relevantes para a dieta]. Gostaria de [inserir o prazo para atingir o objetivo]. Segue algumas informações sobre a minha alimentação atual: [inserir a alimentação diária, detalhando o que come e bebe em cada refeição]".

3. Destaque de restrições alimentares \\
"Crie uma rotina alimentar semanal segura para uma pessoa [inserir a sua restrição alimentar]. Aqui estão as informações que você precisa: [Em formato de lista, informe objetivo, preferências alimentares, números de refeições por dia, condições de preparo, medicação e/ou suplementação (caso haja) e as informações sobre sexo, peso, altura e nível de atividade física que realiza]".

4. Adaptação com a rotina \\
"Crie uma sugestão de alimentos que posso preparar e consumir em até [inserir tempo] durante [inserir período]. Gostaria que esses alimentos fossem [inserir informações sobre o tipo de dieta que busca e contexto sobre o seu pedido, Ex: Trabalho em home office e tenho uma hora de almoço todos os dias]".

5. Orçamento e frequência de compras \\
Monte uma dieta saudável para mim e adequada ao meu orçamento. Faço compras [inserir frequência] e posso gastar até [inserir orçamento]. [Complete o prompt com preferências alimentares, itens que você não consome e outros fatores que possam fazer sentido para a dieta desejada]. 

\section*{Finalmente}
\addcontentsline{toc}{section}{termino}
Não gostou de nenhum desses prompts, transforme a própria IA em um gerador:
\begin{enumerate}
	\item Abra o ChatGPT, Claude, Deepseek ou Gemini.
	\item Digite o seguinte comando: "Você é um especialista em engenharia de prompts. Vou apresentar uma ideia ou descrever o que desejo criar, e você me fará perguntas para entender melhor meus objetivos, esclarecer os requisitos e refinar a proposta. Com base nas minhas respostas, você desenvolverá um prompt claro, completo e bem estruturado, buscando o melhor resultado possível."
	\item Responda às perguntas e copie o resultado gerado.
\end{enumerate}

Sou um entusiasta do mundo \textbf{Open Source} e novas tecnologias. Qual a diferença entre Livre e Open Source? \underline{Livre} significa que esta apostila é gratuita e pode ser compartilhada a vontade. \underline{Open Source} além de livre todos os arquivos que permitem a geração desta (chamados de arquivos fontes) devem ser disponibilizados para que qualquer pessoa possa modificar ao seu prazer, gerar novas, complementar ou fazer o que quiser. Os fontes da apostila (que foi produzida com o LaTex) está disponibilizado no GitHub \cite{github}. Veja ainda outros artigos que publico sobre tecnologia através do meu Blog Oficial \cite{fernandoanselmo}.

%-----------------------------------------------------------------------------
% REFERÊNCIAS
%-----------------------------------------------------------------------------
\begin{thebibliography}{3}
		\bibitem{fernandoanselmo} 
	Fernando Anselmo - Blog Oficial de Tecnologia \\
	\url{http://www.fernandoanselmo.blogspot.com.br/}
	
	\bibitem{publicacao} 
	Encontre essa e outras publicações em \\
	\url{https://cetrex.academia.edu/FernandoAnselmo}
	
	\bibitem{github} 
	Repositório para os fontes da apostila \\
	\url{https://github.com/fernandoans/publicacoes}
\end{thebibliography}

\end{document}